% Discussion: applying the deployability test of Eq. (7). Numbers are macros
% from tables/generated/numbers.tex; the table is generated (EXP-052).
We instantiate Eq.~\eqref{eq:deploy} with bounds an operator can defend. The generalization bound is $\delta = 0.10$ of balanced accuracy, applied to the upper end of the corrected 95\% interval. The alert term asks for a threshold that detects at least half of the attacks ($r_{\min} = 0.5$), makes at least one alert in ten a true one ($\rho = 0.1$ at $\pi = 0.002$), and raises at most $A_{\max} = 200$ false alerts per hour at $\lambda_{b} = 240{,}000$ benign flows per hour, the volume a tier-one team of three analysts can triage after automated filtering. All alert conditions are evaluated on \DB{}, because an operator deploys on the target, not on the source. The latency term uses the strictest budget the specification allows, $B = 10$\,ms: the upper bound of the 99th percentile of the radio decision path meets 1\,ms for every architecture (Section~\ref{sec:res:latency}), so the term passes anywhere in the specified range, with the caveat that indication decoding, control encoding and queueing were not measured.

Table~\ref{tab:predicate} gives the result. No architecture passes: \PredPass{} of six. Every architecture fails the generalization term, with upper bounds of \PredGenMin{} to \PredGenMax{}, and every architecture fails the alert term by a wide margin: at recall of at least 0.5, the best operational precision on \DB{} is \PredPPVmax{} and the fewest false alerts per hour is \PredFAmin{}. The verdict is therefore insensitive to the bounds chosen: it holds for any $\rho$ above \PredPPVmax{}, any $A_{\max}$ below \PredFAmin{} per hour and any $\delta$ below \PredGenMin{}, and it does not depend on the latency term at all. It also survives the label conflicts of Section~\ref{sec:res:target}: evaluated on \DB{} without the benign copies, \PredCleanPass{} of six pass, the best operational precision at recall of at least 0.5 is \PredCleanPPVmax{}, and the fewest false alerts per hour are \PredCleanFAmin. That term matters only once the first two pass; on our emulated measurements it would not decide the verdict.

\begin{table}[t]
\caption{Deployability test of Eq.~\eqref{eq:deploy} per architecture. $\overline{\Delta}_{\mathrm{BA}}$: upper 95\% bound of the balanced-accuracy loss (pass if $\le 0.10$). Best PPV and minimum false alerts per hour over thresholds with recall $\ge 0.5$ on \DB{} (pass if PPV $\ge 0.1$ and alerts $\le 200$). $B_{\min}$: smallest budget met by the upper bound of the radio decision path's 99th percentile, emulated.}
\label{tab:predicate}
\centering
\scriptsize
\setlength{\tabcolsep}{2.4pt}
% GENERATED by analysis/make_tables.py on 2026-09-25
% source: EXP-052 predicate (EXP-041, EXP-043)
% commit: 72269b8
% DO NOT EDIT. Edit the generator, then re-run `make tables`.
\begin{tabular}{lccccccc}
\toprule
Model & $\overline{\Delta}_{\mathrm{BA}}$ & Gen. & Best PPV & Min.\ alerts/h & Alert & $B_{\min}$ (ms) & $\mathcal{P}$ \\
\midrule
LR & 0.32 & $\times$ & 0.0024 & 166\,328 & $\times$ & 1 & 0 \\
DT & 0.50 & $\times$ & 0.0022 & 118\,342 & $\times$ & 1 & 0 \\
RF & 0.47 & $\times$ & 0.0027 & 98\,410 & $\times$ & 1 & 0 \\
XGB & 0.47 & $\times$ & 0.0026 & 95\,507 & $\times$ & 1 & 0 \\
HGB & 0.47 & $\times$ & 0.0026 & 99\,277 & $\times$ & 1 & 0 \\
MLP & 0.50 & $\times$ & 0.0026 & 105\,744 & $\times$ & 1 & 0 \\
\bottomrule
\end{tabular}

\end{table}

\emph{Does the test discriminate?} A test that no detector passes could simply be too strict. The published 5G-NIDD pipeline of Section~\ref{sec:res:target} shows that it is not. On its own random split, with the test's alert term evaluated on that split, \PubRzPassN{} of its five classifiers pass: operational precision reaches \PubRzPPVmax{} at $\pi = 0.002$ with as few as \PubRzFAmin{} false alerts per hour at $\lambda_b = 240{,}000$. The same pipeline fails the test one rung later, when the two position counters are removed, and at every later rung, on all flows and on the flows without the benign copies (Table~\ref{tab:ladder}). The test separates an evaluation that predicts deployment from one that does not; the detectors of this paper fail it because their evaluations are of the second kind.